% !TEX TS-program = pdflatex

\documentclass[12pt]{article}

\usepackage[margin=1.2in]{geometry}
\usepackage{amsmath, amssymb}
\usepackage{booktabs}
\usepackage{listings}
\usepackage{xcolor}
\usepackage{hyperref}
\usepackage{parskip}
\usepackage{microtype}
\usepackage{titlesec}
\usepackage{lmodern}
\usepackage[T1]{fontenc}

% Section styling
\titleformat{\section}{\large\bfseries}{}{0em}{}[\titlerule]
\titleformat{\subsection}{\normalsize\bfseries}{}{0em}{}

% Code block styling
\definecolor{codebg}{RGB}{248,248,248}
\definecolor{codecomment}{RGB}{100,100,100}
\definecolor{codekeyword}{RGB}{0,0,180}
\definecolor{codestring}{RGB}{160,0,0}

\lstset{
  backgroundcolor=\color{codebg},
  basicstyle=\ttfamily\small,
  commentstyle=\color{codecomment}\itshape,
  keywordstyle=\color{codekeyword}\bfseries,
  stringstyle=\color{codestring},
  breaklines=true,
  frame=single,
  framesep=6pt,
  rulecolor=\color{gray!40},
  xleftmargin=12pt,
  xrightmargin=12pt,
  aboveskip=10pt,
  belowskip=10pt,
  showstringspaces=false
}

\hypersetup{
  colorlinks=true,
  linkcolor=black,
  urlcolor=blue!70!black
}

% -----------------------------------------------
\title{%
  \vspace{-1cm}
  \LARGE\bfseries The Future of Programming:\\[0.4em]
  \large From Problem Specifications to AI-Compiled Implementations ---\\
  A Type-Theoretic Foundation
}
\author{Ashwin Rao\\
  \small 
  Adjunct Professor, ICME, Stanford University
}
\date{\today}
% -----------------------------------------------

\begin{document}

\maketitle
\vspace{0.5em}
\hrule
\vspace{1.2em}

\begin{abstract}
The software industry has spent seventy years getting progressively better at
describing \emph{how} to implement solutions, while remaining remarkably
imprecise about describing \emph{what problem} is actually being solved. This
essay argues that we are on the verge of a fundamental correction: driven by
the capabilities of modern AI, the act of programming will shift from writing
implementation code to writing formal, structured specifications of problems.
We trace the historical and theoretical roots of this shift --- from Church's
lambda calculus and the Curry--Howard correspondence to modern dependent type
theory --- and propose a realistic two-stage compilation pipeline in which AI
transforms human problem descriptions into formal typed specifications, and
then into verified implementations. The intended audience is anyone with a
working familiarity with computer science and modern AI.
\end{abstract}

\vspace{0.8em}

% ===============================================
\section{The Future of Programming}
% ===============================================

We are on the verge of a more fundamental shift in programming than most
people in the industry currently appreciate. Not just better tools, faster
code generation, or smarter autocomplete --- but a change in what programming
\emph{is}.

For most of the history of software, brilliant engineers have dived into
\emph{how} to build something before rigorously nailing down \emph{what
problem} they were actually solving. This is not simply a matter of
professional discipline. It is a deeper, structural problem: a seventy-year
architectural choice baked into how the industry was built. Understanding why
requires a short detour into the theoretical foundations of computing.

% ===============================================
\section{Two Formalisms, One Fork in the Road}
% ===============================================

In the 1930s, two foundational models of computation emerged almost
simultaneously, and they encoded fundamentally different habits of mind.

Alan Turing's machine described computation as a sequence of mechanical state
transitions --- an irreducibly \emph{how}-oriented formalism. To specify a
computation in Turing's framework is to describe, in meticulous low-level
detail, every step the machine must take.

Alonzo Church's lambda calculus described computation as the composition of
mathematical functions --- a \emph{what}-oriented formalism. To specify a
computation in Church's framework is to define what the function \emph{is}:
its inputs, its outputs, and the precise relationship between them.

The Church--Turing thesis established that these two formalisms are
computationally equivalent --- anything expressible in one is expressible in
the other. But this equivalence conceals a crucial asymmetry in epistemic
style. \textbf{In Church's framework, the solution specification and the
problem specification are inseparable.} You cannot write a well-typed,
compositional function without having already characterised precisely what
problem it solves. To define the domain and codomain of a function is to
formally state what you are computing and why. Turing's approach carries no
such obligation: one can write a sequence of machine operations without ever
articulating the problem being addressed.

Functional programming languages --- Haskell being the clearest modern
exemplar --- carry Church's spirit: they encourage the programmer to think
first in terms of \emph{what} a computation is before considering \emph{how}
it will execute.

% ===============================================
\section{How the Industry Went the Other Way}
% ===============================================

Despite the theoretical elegance of Church's approach, it was Turing's model
that the industry adopted, and for understandable reasons. Early computers
were, in essence, physical Turing machines: sequential, imperative,
memory-constrained devices. Imperative languages mapped directly onto that
hardware and ran efficiently. Functional languages, being closer to
mathematical abstractions, required more sophisticated compilers and more
memory --- both scarce resources in the early decades of computing. Once C
dominated systems programming, the entire ecosystem of tooling, education,
and hiring calcified around the imperative model.

The resulting historical arc is familiar to every programmer: Turing
machines, then punched cards, then assembly, then C, then C++, then Java,
then Python. Each generation raised the level of abstraction over
\emph{implementation details}, making programming progressively more
accessible. But throughout this entire arc, the industry never fundamentally
crossed over to Church's side of the ledger. Each step made it easier to
describe solutions. None of them made it more rigorous to describe problems.

The consequences have been severe. When implementation mechanics take
priority over problem specification, the result is poorly designed systems,
unnecessary complexity, and an ever-growing mountain of technical debt. A
very large fraction of the complexity in modern software systems is
\emph{accidental} rather than \emph{essential} --- it arises not from the
inherent difficulty of the problem but from the failure to understand and
specify the problem rigorously before reaching for a solution.

% ===============================================
\section{Formal Methods: The Road Not Taken}
% ===============================================

It is worth noting that the idea of formally specifying problems before
implementing solutions is not new. A research tradition of formal
specification languages has existed for decades. Tools such as Z~notation,
VDM, TLA+, Alloy, Coq, and Agda all represent attempts to bring
mathematical rigour to the problem specification side of software development.

These efforts never went mainstream. The reason is instructive: even when a
formal specification existed, a human expert still had to bridge the gap
between the spec and a working implementation. That gap demanded rare
expertise and enormous effort, and so was never crossed at the scale the
industry required. Formal methods remained a specialised discipline, confined
largely to safety-critical domains such as aerospace and cryptographic
protocols.

This history raises an important question: if formal specification languages
have existed for decades without achieving adoption, what is different now?
The answer is AI.

% ===============================================
\section{Why AI Changes Everything}
% ===============================================

AI can serve as the compiler between a precise problem specification and a
working implementation --- exactly as high-level language compilers once
spared programmers from writing assembly code. The human's job is to specify
\emph{what}. The machine's job is to produce \emph{how}. The gap that
defeated formal methods for decades --- the translation from spec to code ---
is precisely the kind of task at which modern large language models, guided
by formal type systems, are increasingly capable.

This reframes the future of programming entirely. The next generation of
programming will not be about writing better code. It will be about writing
better \emph{problem specifications}. We will converge on a formal language
purpose-built for exactly that: a structured, precise language for describing
what you want solved, which AI then compiles into a verified implementation.
\textbf{This problem specification language will be the new programming language
and the problem specifications we write in this language will be the new source code}.

The remainder of this paper addresses two questions. First: what is the right
mathematical foundation for such a specification language? Second: given that
most practitioners cannot be expected to write mathematical specifications
from scratch, what does a realistic and deployable architecture look like?

This requires a bit of a detour into mathematics, and specifically into \emph{Type Theory}.

% ===============================================
\section{What Is Type Theory?}
% ===============================================

Type theory, originating with Bertrand Russell and substantially developed by
Per Martin-L\"{o}f in the 1970s, is a foundational system for mathematics and
logic in which every object carries a \emph{type}, and types are themselves
mathematical objects that can be reasoned about. The central intuition is
simple but profound:

\begin{center}
\textbf{A type is a proposition. A program that inhabits that type is a proof
of that proposition.}
\end{center}

This equivalence is known as the \textbf{Curry--Howard correspondence}, and it
is one of the deepest results in theoretical computer science. It establishes
a three-way isomorphism between logic, type theory, and programming:

\vspace{0.6em}
\begin{center}
\begin{tabular}{@{}lll@{}}
\toprule
\textbf{Logic} & \textbf{Type Theory} & \textbf{Programming} \\
\midrule
Proposition      & Type                 & Specification         \\
Proof            & Term                 & Program               \\
True proposition & Inhabited type       & Working implementation \\
\bottomrule
\end{tabular}
\end{center}
\vspace{0.6em}

This is not metaphor --- it is mathematically precise. When you write a type,
you state a claim about the world. When you write a program of that type, you
prove that claim constructively. The direct lineage from Church's lambda
calculus is clear: Church said a computation is a function; type theory says
the function's type is a proposition about what that function computes; a
program is the proof that the proposition is satisfiable.

% ===============================================
\section{A Spectrum of Expressiveness}
% ===============================================

Most programmers encounter type theory in diluted form through the type
systems of their languages. There is a spectrum of expressiveness worth
understanding.

\subsection{Simple types}

In languages such as C and Java, types classify data: \texttt{int},
\texttt{string}, \texttt{bool}. They catch basic errors but say almost
nothing about behaviour. The signature

\begin{lstlisting}[language=Java]
int add(int x, int y)
\end{lstlisting}

tells us the function accepts two integers and returns one. It says nothing
about what the function actually computes.

\subsection{Parametric and generic types}

Languages such as Haskell and Java (with generics) allow types to be
quantified over type variables: \texttt{List<T>}, \texttt{Functor f}. This
enables structural correctness guarantees through abstraction --- a
\texttt{map} function that is polymorphic over \texttt{List} cannot, by
construction, drop or reorder elements.

\subsection{Dependent types}

Dependent type systems --- found in Coq, Agda, Lean, and Idris --- are where
type theory becomes a genuine specification language. In a dependently typed
system, \textbf{types can depend on values}. This allows arbitrary
propositions to be encoded as types.

Consider a length-indexed vector type:
\[
  \texttt{Vector} : \mathbb{N} \to \texttt{Type} \to \texttt{Type}
\]
A value of type $\texttt{Vector}\; n\; a$ is a list of \emph{exactly} $n$
elements of type $a$. The function

\begin{lstlisting}[mathescape=true]
append : Vector m a $\to$ Vector n a $\to$ Vector (m+n) a
\end{lstlisting}

is not merely typed --- it is \emph{proven correct by construction}. An
implementation returning a vector of the wrong length will not typecheck. The
compiler is the verifier.

Further examples of propositions expressible as dependent types:

\begin{itemize}
  \item \texttt{SortedList a} --- a list guaranteed by its type to be sorted;
        no runtime assertion required.
  \item \texttt{\{ x : Int | x > 0 \}} --- a refinement type encoding a
        positive integer; the constraint is in the type, not buried in an
        assertion.
  \item \texttt{Authenticated (Request r)} --- a request type that can only
        be constructed after authentication has been verified; security as a
        type invariant.
\end{itemize}

In a fully dependently typed language, the distinction between \emph{writing
a specification} and \emph{writing a proof} dissolves entirely. The type
\emph{is} the specification; the program \emph{is} the evidence that the
specification is satisfiable.

% ===============================================
\section{Type Theory as the Foundation of the Problem Spec Language}
% ===============================================

We can now make the connection to our central thesis precise. Dependent type
theory provides the natural mathematical substrate for the future problem
specification language. A problem specification in this framework might look
as follows:

\begin{lstlisting}[mathescape=true]
-- Specify the problem: sort a list
sort : (a : Type) $\to$ Ord a $\to$ (input : List a)
     $\to$ { result : List a | sorted result $\wedge$ permutation result input }
\end{lstlisting}

This type asserts: given any ordered type \texttt{a} and an input list,
produce a result that is \emph{both} sorted \emph{and} a permutation of the
input. This is not a comment, a test, or a docstring --- it is a mathematical
proposition. Any program inhabiting this type is, by construction, a correct
sorting algorithm. The typechecker does not probabilistically check
correctness; it verifies it.

Non-functional requirements (NFR) --- latency, throughput, security, compliance ---
are first-class citizens of this framework. Refinement types and dependent
types encode them as typed constraints directly in the specification:

\begin{lstlisting}[mathescape=true]
processPayment
  : (req : PaymentRequest)
  $\to$ Latency p95 < 50ms
  $\to$ Compliant req PCI_DSS
  $\to$ { result : PaymentResult | settled result $\vee$ explicitly_failed result }
\end{lstlisting}

The same functional specification, composed with different non-functional
constraint types, produces a different valid implementation space --- and the
AI compiler navigates that space accordingly. This directly addresses one of
the most persistent sources of late-discovered complexity in software
projects: non-functional requirements (NFR) that live informally in design documents
and tribal knowledge rather than as first-class, verifiable constraints.

% ===============================================
\section{A Realistic Two-Stage Compilation Pipeline}
% ===============================================

There is an important practical objection to the vision so far. Dependent
type theory is formidably inaccessible to most practitioners. Expecting
engineers, domain experts, and architects to write specifications in Coq or
Agda is not a realistic near-term proposition --- and it would recreate the
same accessibility problem that kept formal methods from going mainstream in
the first place.

The resolution is that the right architecture is not a single compilation
step from human intent to implementation, but \textbf{two}, with a formal
typed specification serving as a crucial intermediate representation.

\subsection{Stage 1: The human problem specification}

The starting point is a problem description authored by a human in whatever
combination of formalisms they are comfortable with. This is not a free-form
chat prompt --- it is a deliberate, structured artefact, composed with as
much clarity and precision as the author can bring. It might include:

\begin{itemize}
  \item \textbf{Natural language}, written with precision, structure, and
        detail --- describing goals, constraints, edge cases, and intent;
  \item \textbf{Existing specification documents} --- prior design documents,
        architecture decision records, regulatory requirements, or
        organisational standards that can be referenced or attached;
  \item \textbf{References to current systems} --- existing data schemas,
        APIs, software interfaces, and data assets that the solution must
        interact with or conform to;
  \item \textbf{Semi-formal artefacts} for more technical authors: JSON or
        YAML configurations, typed schemas, mathematical equations,
        ontologies, functional specifications, and structured descriptions
        of non-functional requirements (NFR) such as latency budgets, throughput
        targets, security constraints, and compliance obligations.
\end{itemize}

The key principle is that \textbf{more structure always helps}. A natural
language description is a valid starting point; the same description
augmented with typed schemas, NFR budgets, and references to existing systems
is a substantially better one. The human is not expected to achieve formal
precision --- that is the machine's responsibility in the next stage --- but
they are expected to bring as much clarity and structure as they can.

\subsection{Stage 2: AI compilation to a formal dependent-typed specification}

The first AI compilation step transforms the human problem specification into
a formal, dependent-typed machine specification. This is not a trivial
translation. It requires the AI to resolve ambiguities, infer missing
constraints, and make the implicit explicit. The output is a dependent-typed
specification that is:

\begin{itemize}
  \item \textbf{Precise}: every input, output, invariant, and constraint is
        formally stated;
  \item \textbf{Verifiable}: the typechecker can confirm internal
        consistency;
  \item \textbf{Inspectable}: a human engineer can read, audit, and edit it
        before compilation proceeds;
  \item \textbf{Machine-actionable}: it provides an unambiguous target for
        the subsequent implementation synthesis step.
\end{itemize}

This intermediate representation is the linchpin of the entire system.
Without it, there is no principled mechanism to verify that the human's
intent has been correctly understood. With it, correctness ceases to be a
matter of interpretation --- it becomes a mathematical property.

The AI iterates with the human at this stage to confirm that the typed
specification faithfully captures the intended problem. Ambiguities are
surfaced and resolved in dialogue before any implementation is generated.
This iteration loop between human intent and machine formalisation is where
much of the real value of the pipeline lies.

\subsection{Stage 3: AI compilation to implementation}

The second AI compilation step takes the formal dependent-typed specification
and synthesises a working implementation --- a program that provably inhabits
the specified type. This is type-directed program synthesis: the AI searches
the space of programs for one whose type matches the specification, and the
typechecker verifies the result with mathematical certainty.

The full pipeline is therefore:

\vspace{0.8em}
\begin{center}
\begin{tabular}{@{}ccccc@{}}
\toprule
\textbf{Human spec} & $\xrightarrow{\;\text{AI}\;}$ &
\textbf{Typed machine spec} & $\xrightarrow{\;\text{AI}\;}$ &
\textbf{Implementation} \\
\midrule
Natural language, & & Dependent types, & & Verified program \\
semi-formal artefacts, & & NFR constraints, & & inhabiting the \\
system references & & formal invariants & & specified type \\
\bottomrule
\end{tabular}
\end{center}
\vspace{0.8em}

At each stage, \textbf{human intervention is optional but meaningful}. A
technically sophisticated engineer may inspect and edit the dependent-typed
machine specification --- tightening constraints, adding invariants, or
correcting a misunderstood requirement --- before compilation proceeds to
implementation. An engineer may similarly inspect and edit the generated
implementation. Crucially, \textbf{every edit is checked for consistency
against the dependent-typed specification}: the typechecker enforces that no
manual modification can introduce a violation of the formally stated
requirements. Correctness integrity is maintained throughout.

\subsection{Why the intermediate specification is essential}

One might ask: why not compile directly from the human problem description to
the implementation, bypassing the intermediate typed specification? The answer
is that without the formal intermediate representation, the system has no
principled mechanism to verify that it has correctly understood the human's
intent. Natural language is ambiguous; informal artefacts are incomplete;
references to existing systems carry implicit assumptions. A direct
compilation to code produces something that may look correct but carries no
formal guarantee.

The dependent-typed intermediate specification makes the AI's interpretation
of the human's intent \emph{explicit and auditable}. It is the point in the
pipeline where intent becomes verifiable fact. Removing it would be analogous
to removing the type system from a programming language: the code might still
run, but the guarantees disappear.

\subsection{Why the human specification is the right starting point}

Conversely, one might ask: why not require all practitioners to write
dependent types directly, eliminating the first compilation step? The answer
is that the population of people who possess essential problem knowledge is
far larger than the population of type theorists. Domain experts, product
engineers, architects, and data scientists all carry knowledge that is
indispensable to a correct specification. Requiring type-theoretic literacy
as a precondition would recreate the accessibility barrier that kept formal
methods from achieving adoption for decades.

The two-stage pipeline democratises access to formal verification. The human
contributes \emph{knowledge and intent}; the machine contributes
\emph{precision and proof}. Neither can substitute for the other.

% ===============================================
\section{The Emerging Stack}
% ===============================================

Assembling these pieces, we can sketch the future programming stack clearly:

\begin{enumerate}
  \item \textbf{Human writes a structured problem specification}, in natural
        language augmented with whatever semi-formal artefacts they can
        bring --- typed schemas, NFR constraints, references to existing
        systems and data. Clarity and structure are the primary virtues at
        this stage.

  \item \textbf{AI compiles the human specification into a dependent-typed
        machine specification}, iterating with the human to resolve
        ambiguities and confirm that intent has been correctly formalised.
        This is the new ``source of truth'' for the system.

  \item \textbf{AI synthesises a verified implementation} from the typed
        specification --- a program that provably inhabits the specified
        type. The typechecker verifies correctness with mathematical
        certainty, replacing the testing and debugging loop for all
        properties expressible in the specification.

  \item \textbf{Humans optionally inspect and edit} either the typed
        specification or the implementation, with the typechecker enforcing
        consistency at every step.
\end{enumerate}

This is already a live research direction. Program synthesis from types,
type-directed code generation using large language models, and AI-assisted
formal verification are all active areas at the intersection of programming
languages theory and machine learning.

% ===============================================
\section{Why Type Theory Is the Right Foundation}
% ===============================================

Three properties make dependent type theory the natural mathematical home for
the problem specification language, rather than natural language or informal
diagrams alone.

\textbf{Compositionality.} Types compose. A complex system specification is
built from smaller typed components, each independently verifiable, in exactly
the way Church envisioned functions composing. This is the direct lineage from
lambda calculus to the future of programming.

\textbf{NFR expressiveness.} Refinement types and dependent types make
non-functional requirements --- latency bounds, memory limits, security
invariants, compliance properties --- first-class constraints in the
specification rather than afterthoughts in a Confluence page. They are
verified by the typechecker with the same rigour as functional correctness.

\textbf{AI alignment.} A dependently typed specification gives a coding agent
an unambiguous, machine-checkable target. The agent is not guessing intent
from a vague prompt; it is searching a well-defined space of programs for one
that satisfies a formal proposition. The typechecker closes the loop with
certainty rather than probability.

% ===============================================
\section{Conclusion}
% ===============================================

The through-line from Church to here is remarkably clean.

Church said a computation is a function. Martin-L\"{o}f said a function's
type is a proposition. Curry and Howard said a program is a proof. The vision
articulated in this paper says: the human states the proposition, the AI
formalises it, and the AI finds the proof.

The industry spent seventy years getting better at the \emph{how}. The next
era will be defined by getting rigorous about the \emph{what}. The ``new
programming language'' is, in its mathematical essence, a dependently typed
specification language --- one that most humans will approach through natural
language and semi-formal artefacts, and that AI will compile first into formal
precision and then into verified code.

The future of programming is not less rigorous. It is more rigorously focused
on the right thing.

\end{document}